\documentclass[11pt]{article}
\usepackage[margin=2.1cm]{geometry}
\usepackage{booktabs}
\usepackage{amsmath}
\usepackage{enumitem}
\usepackage{xcolor}
\usepackage[T1]{fontenc}
\usepackage{lmodern}

\setlength{\parindent}{0pt}
\setlength{\parskip}{3.6pt}
\newcommand{\code}[1]{\texttt{#1}}
\newcommand{\hr}{\vspace{2pt}\hrule\vspace{6pt}}

\begin{document}

\begin{center}
{\Large\bfseries EE\,362 --- Computer Organization Laboratory}\\[2pt]
{\large\bfseries Lab 12 --- CPU Integration}\\[4pt]
Fall 2026 \quad$\cdot$\quad Lab TA: Gandhi Torres \quad$\cdot$\quad Sections: Mon / Tue, 3:45--5:45~PM, ECIII 237
\end{center}
\hr

\section*{1. Objective}

Wire twelve components together and run a program.

\begin{quote}
\emph{Debugging a system means predicting it first.}
\end{quote}

\textbf{Almost no new VHDL.} This lab is wiring, timing and debugging. Expect it
to be emotionally hard and technically shallow: nothing here is difficult to
write, and finding out why it does not work is the entire exercise.

\textbf{Bring your printed control table.} Step 3 of the debugging method needs
it.

\section*{2. What success looks like}

\begin{center}
\begin{tabular}{ll}
\toprule
instructions executed & \textbf{45}, one per clock period, CPI $= 1$ \\
final \code{r1} & \code{0x00000000} \quad the loop counter ran out \\
final \code{r2} & \code{0x00000400} \quad $1024 = 2^{10}$ \\
word stored at address 0 & \code{0x00000400} \\
\bottomrule
\end{tabular}
\end{center}

\textbf{Your CPU computes $2^{10}$.} Every component you have built since Lab 1
exists to make that happen.

\section*{3. The timing contract}

Write this on paper before you wire anything.

\begin{center}
\begin{tabular}{ll}
\toprule
\textbf{rising edge} & \textbf{FETCH.} \quad \code{IRreg <= ROM[PC]} \\
\textbf{falling edge} & \textbf{COMMIT.} \quad register file writes, RAM writes, PC updates \\
\bottomrule
\end{tabular}
\end{center}

One instruction per clock period. The instruction is captured at the start of
the period; every piece of state changes in the middle of it, by which time the
combinational logic has settled.

\textbf{The register file and the NIAG need an inverted clock.} Their own
architectures use \code{rising\_edge}, and you want them to commit on the
falling edge with everything else. Inverting the clock at the instantiation is
how you retime a component you do not want to modify.

\textbf{Every bug in this lab that is not a mis-wire is an edge mistake.}

\section*{4. Three wiring decisions that are easy to get wrong}

\begin{enumerate}[leftmargin=1.8em,itemsep=3pt]
\item \textbf{\code{Asel = rt} and \code{Bsel = rs}.} Not the other way round.
Work it out rather than trusting the handout:

\begin{center}
\begin{tabular}{ll}
\toprule
\code{ADDI}/\code{LW}/\code{SW} & the base register is \code{rs} and must arrive on \code{Bout}, \\
 & because the ALU's A input is the immediate \\
\code{SLL} & shifts \code{rt}, which must arrive on \code{Aout} \\
\code{SW} & stores \code{rt}, driven onto the bus from \code{Aout} \\
\code{BLEZ} & has \code{rt} $= 0$, so \code{Aout + Bout} $= 0 + rs$ \\
\bottomrule
\end{tabular}
\end{center}

One assignment satisfies all eight instructions. \textbf{Swap them and three of
the four instruction types still work} --- which is what makes it nasty.

\item \textbf{The store value comes from \code{Aout}}, through the
\code{TriBuffer}.

\item \textbf{The ROM is addressed by \code{PCreg}, not by \code{NSA}.} Address
it with \code{NSA} and every instruction executes one too early.
\end{enumerate}

\section*{5. The debugging method}

When it does not work --- and it will not, the first time --- \textbf{do not
open the waveform and stare at forty signals}. Bisect. Four checks, in this
order, every time.

\begin{center}
\begin{tabular}{cll}
\toprule
& check & if it fails, the fault is in \\
\midrule
1 & Is the PC advancing by 4? & the NIAG, or reset \\
2 & Is \code{IR} the instruction the ROM table says? & ROM addressing, or the fetch process \\
3 & Do the control signals match your table? & the decoder \\
4 & Only now, the datapath. & the wiring between components \\
\bottomrule
\end{tabular}
\end{center}

Each check depends on the ones before it. Checking the ALU while the PC is stuck
tells you nothing --- the ALU is faithfully computing on garbage.

\textbf{Do them in order even when you are sure you know what is wrong}, because
you are usually not. Example 1 gives you a CPU with two deliberate wiring errors
to practise on; both pass checks 1, 2 and 3, and establishing that is not wasted
effort --- it is what tells you the fault is in the wiring rather than the
control.

\section*{6. Hand-compute before you simulate}

\textbf{Work out the first five rows of the trace on paper first.} This is
graded, and not only because it is good practice: a student who only ever reads
values off a waveform has no way to tell a wrong trace from a right one.

\begin{center}
\begin{tabular}{cllll}
\toprule
cyc & PC & instruction & what it does & register written \\
\midrule
0 & \code{0x00} & \code{LUI r1, 11} & $r_1 \leftarrow 11$ & \code{r1} $=$ \code{0x0000000B} \\
1 & \code{0x04} & \code{LUI r2, 1} & $r_2 \leftarrow 1$ & \code{r2} $=$ \code{0x00000001} \\
2 & \code{0x08} & \code{ADDI r1, r1, -1} & $r_1 \leftarrow 10$ & \code{r1} $=$ \code{0x0000000A} \\
3 & \code{0x0C} & \code{BLEZ r1, +2} & $10 > 0$, not taken & --- \\
4 & \code{0x10} & \code{SLL r2, r2, 1} & $r_2 \leftarrow 2$ & \code{r2} $=$ \code{0x00000002} \\
\bottomrule
\end{tabular}
\end{center}

\section*{7. Assignment}

\textbf{Step 1:} replace anything in \code{components/} with your own version.
The reference copies exist so that nobody is blocked by an earlier lab.
\textbf{Using a reference component costs you nothing on this lab} --- the
earlier lab's grade simply stands as it was. Integration is never gated on an
earlier miss; that is a deliberate rule of this course, not a favour.

\textbf{Step 2:} complete \code{CPU.vhd}. Twelve instantiations, two clocked
processes, named association throughout.

The testbench is \textbf{provided complete} --- this lab is wiring and
debugging, not test writing. It prints the execution trace and checks the final
state.

\subsection*{One thing in the testbench worth reading}

The tracer samples \emph{just before} the falling edge, never after it. After
the edge the write has committed, and for a self-referencing instruction such as
\code{ADDI r1, r1, -1} the combinational logic immediately recomputes from the
\emph{new} \code{r1} --- so a post-edge sample would show the result of the next
iteration. The hardware would be correct and the observer would be lying.

\textbf{When you sample is part of the measurement.} This is a real and common
way to waste an hour.

\subsection*{On the metavalue warnings}

You will see several \code{NUMERIC\_STD.TO\_INTEGER: metavalue detected}
warnings in the first few cycles. They are normal: before the first instruction
is fetched some addresses are still \code{'U'}. They stop once the machine is
running --- and a metavalue warning that \emph{never} stops means a signal is
never being driven, which makes them the first thing to look at if your CPU is
dead.

\section*{8. Evidence to collect --- E1 through E7, all of them}

This is the only design all semester that reports every evidence source
meaningfully.

\begin{center}
\begin{tabular}{lll}
\toprule
& Source & What it must show \\
\midrule
E1 & Flow Summary + Messages & 0 errors; warnings sorted \\
E2 & RTL Viewer & the twelve components, wired \\
E3 & Technology Map Viewer & what the fabric made of it \\
E4 & Resource Utilisation & PC 32, IR 64, register file $\sim$1024, block memory \\
E5 & Pin-Out & \textbf{one} functional pin, \code{SysClk} \\
E6 & Timing Analyzer & Fmax, and the critical path \\
E7 & Waveform + transcript & $\sim$20 annotated cycles, and \code{PASS} \\
\bottomrule
\end{tabular}
\end{center}

\textbf{E5 is worth a sentence.} In Lab 5 the ALU needed 108 pins as a top-level
entity. Here the whole processor needs one clock. That contrast is what ``wide
blocks live inside the chip'' means.

\section*{9. Report --- \code{lab12.tex}, the longest chapter}

\begin{itemize}[leftmargin=1.6em,itemsep=1pt]
\item the timing contract, in your own words;
\item the \textbf{full 45-row execution trace};
\item your \textbf{hand-computed first five rows}, beside the simulated ones;
\item an annotated waveform of $\sim$20 cycles;
\item the complete evidence set E1--E7, each with a sentence;
\item \textbf{Discussion:} the integration bugs you hit, and which of the four
checks found each one.
\end{itemize}

Expect 8--12 pages. This is the chapter that makes the report a portfolio piece.

\section*{10. Grading (100 points)}

\begin{center}
\begin{tabular}{lr}
\toprule
Pre-lab, submitted at the start of the session & 5 \\
Design compiles cleanly, warnings explained & 25 \\
Hand-computed trace and verification & 20 \\
Simulation results: the CPU runs the program & 20 \\
Report chapter quality & 15 \\
In-lab demonstration and explanation & 15 \\
\midrule
Total & 100 \\
\bottomrule
\end{tabular}
\end{center}

Late work is accepted for one week at $-20$ points.

\section*{11. Things that will go wrong}

\begin{itemize}[leftmargin=1.6em,itemsep=1pt]
\item \textbf{Positional \code{port map}.} Twelve components is where this stops
being survivable.
\item \textbf{Clock-edge mistakes.} The register file and NIAG need the inverted
clock.
\item \textbf{\code{'U'} propagating from one uninitialised signal.} Almost
always the real cause of a dead CPU.
\item \textbf{\code{Asel}/\code{Bsel} swapped.} Three instruction types keep
working.
\item \textbf{\code{TriBuffer} and \code{RAM} both driving the bus.} Look for
\code{'X'}.
\item \textbf{Reading an \code{out} port.} Use an internal signal.
\end{itemize}

\vspace{10pt}
\hr
{\small Three months ago you were writing a four-row truth table for a latch.
Same loop, every week --- specification, truth table, equations, VHDL,
simulation, check. That is how this was always going to end.}

\end{document}
